% Options for packages loaded elsewhere
\PassOptionsToPackage{unicode}{hyperref}
\PassOptionsToPackage{hyphens}{url}
\documentclass[
]{article}
\usepackage{xcolor}
\usepackage[margin=1in]{geometry}
\usepackage{amsmath,amssymb}
\setcounter{secnumdepth}{-\maxdimen} % remove section numbering
\usepackage{iftex}
\ifPDFTeX
  \usepackage[T1]{fontenc}
  \usepackage[utf8]{inputenc}
  \usepackage{textcomp} % provide euro and other symbols
\else % if luatex or xetex
  \usepackage{unicode-math} % this also loads fontspec
  \defaultfontfeatures{Scale=MatchLowercase}
  \defaultfontfeatures[\rmfamily]{Ligatures=TeX,Scale=1}
\fi
\usepackage{lmodern}
\ifPDFTeX\else
  % xetex/luatex font selection
\fi
% Use upquote if available, for straight quotes in verbatim environments
\IfFileExists{upquote.sty}{\usepackage{upquote}}{}
\IfFileExists{microtype.sty}{% use microtype if available
  \usepackage[]{microtype}
  \UseMicrotypeSet[protrusion]{basicmath} % disable protrusion for tt fonts
}{}
\makeatletter
\@ifundefined{KOMAClassName}{% if non-KOMA class
  \IfFileExists{parskip.sty}{%
    \usepackage{parskip}
  }{% else
    \setlength{\parindent}{0pt}
    \setlength{\parskip}{6pt plus 2pt minus 1pt}}
}{% if KOMA class
  \KOMAoptions{parskip=half}}
\makeatother
\usepackage{longtable,booktabs,array}
\newcounter{none} % for unnumbered tables
\usepackage{calc} % for calculating minipage widths
% Correct order of tables after \paragraph or \subparagraph
\usepackage{etoolbox}
\makeatletter
\patchcmd\longtable{\par}{\if@noskipsec\mbox{}\fi\par}{}{}
\makeatother
% Allow footnotes in longtable head/foot
\IfFileExists{footnotehyper.sty}{\usepackage{footnotehyper}}{\usepackage{footnote}}
\makesavenoteenv{longtable}
\usepackage{graphicx}
\makeatletter
\newsavebox\pandoc@box
\newcommand*\pandocbounded[1]{% scales image to fit in text height/width
  \sbox\pandoc@box{#1}%
  \Gscale@div\@tempa{\textheight}{\dimexpr\ht\pandoc@box+\dp\pandoc@box\relax}%
  \Gscale@div\@tempb{\linewidth}{\wd\pandoc@box}%
  \ifdim\@tempb\p@<\@tempa\p@\let\@tempa\@tempb\fi% select the smaller of both
  \ifdim\@tempa\p@<\p@\scalebox{\@tempa}{\usebox\pandoc@box}%
  \else\usebox{\pandoc@box}%
  \fi%
}
% Set default figure placement to htbp
\def\fps@figure{htbp}
\makeatother
\setlength{\emergencystretch}{3em} % prevent overfull lines
\providecommand{\tightlist}{%
  \setlength{\itemsep}{0pt}\setlength{\parskip}{0pt}}
\usepackage{bookmark}
\IfFileExists{xurl.sty}{\usepackage{xurl}}{} % add URL line breaks if available
\urlstyle{same}
\hypersetup{
  hidelinks,
  pdfcreator={LaTeX via pandoc}}

\author{}
\date{\vspace{-2.5em}}

\begin{document}

\section*{Notation}\label{notation}
\addcontentsline{toc}{section}{Notation}

This page lists the symbols, conventions, and abbreviations used
throughout the book. Where a chapter departs from these conventions
locally, the chapter says so explicitly.

\subsection*{Greek symbols}\label{greek-symbols}
\addcontentsline{toc}{subsection}{Greek symbols}

{\def\LTcaptype{none} % do not increment counter
\begin{longtable}[]{@{}
  >{\raggedright\arraybackslash}p{(\linewidth - 2\tabcolsep) * \real{0.2317}}
  >{\raggedright\arraybackslash}p{(\linewidth - 2\tabcolsep) * \real{0.7683}}@{}}
\toprule\noalign{}
\begin{minipage}[b]{\linewidth}\raggedright
Symbol
\end{minipage} & \begin{minipage}[b]{\linewidth}\raggedright
Meaning
\end{minipage} \\
\midrule\noalign{}
\endhead
\bottomrule\noalign{}
\endlastfoot
\(\alpha\) & Type I error rate (significance threshold). \\
\(\beta\) & Type II error rate; also a regression coefficient
(context). \\
\(\mu\) & Population mean. \\
\(\sigma\) & Population standard deviation; \(\sigma^2\) for
variance. \\
\(\pi\) & Population proportion. \\
\(\theta\) & Generic population parameter; \(\hat{\theta}\) for the
estimator. \\
\(\tau^2\) & Between-study variance in meta-analysis. \\
\(\rho\) & Population correlation; intra-class correlation (ICC). \\
\(\lambda\) & Rate parameter; eigenvalue (context). \\
\end{longtable}
}

\subsection*{Latin symbols}\label{latin-symbols}
\addcontentsline{toc}{subsection}{Latin symbols}

{\def\LTcaptype{none} % do not increment counter
\begin{longtable}[]{@{}
  >{\raggedright\arraybackslash}p{(\linewidth - 2\tabcolsep) * \real{0.2317}}
  >{\raggedright\arraybackslash}p{(\linewidth - 2\tabcolsep) * \real{0.7683}}@{}}
\toprule\noalign{}
\begin{minipage}[b]{\linewidth}\raggedright
Symbol
\end{minipage} & \begin{minipage}[b]{\linewidth}\raggedright
Meaning
\end{minipage} \\
\midrule\noalign{}
\endhead
\bottomrule\noalign{}
\endlastfoot
\(n\) & Sample size. \\
\(N\) & Population size. \\
\(\bar{x}\) & Sample mean. \\
\(s\) & Sample standard deviation; \(s^2\) for variance. \\
\(\hat{p}\) & Sample proportion. \\
\(\mathrm{SE}\) & Standard error of an estimator. \\
\(p\) & \(p\)-value. \\
\(L\), \(\ell\) & Likelihood and log-likelihood. \\
\(I^2\) & Heterogeneity statistic in meta-analysis. \\
\end{longtable}
}

\subsection*{Operators and notation}\label{operators-and-notation}
\addcontentsline{toc}{subsection}{Operators and notation}

\begin{itemize}
\tightlist
\item
  \(\sim\) --- ``is distributed as''.
\item
  \(\hat{\theta}\) --- an estimator of the parameter \(\theta\).
\item
  \(X \perp Y\) --- independence.
\item
  \(E[\cdot]\), \(\mathrm{Var}[\cdot]\) --- expectation and variance.
\item
  \(X \mid Y\) --- \(X\) given \(Y\) (conditional).
\end{itemize}

\subsection*{Abbreviations}\label{abbreviations}
\addcontentsline{toc}{subsection}{Abbreviations}

{\def\LTcaptype{none} % do not increment counter
\begin{longtable}[]{@{}
  >{\raggedright\arraybackslash}p{(\linewidth - 2\tabcolsep) * \real{0.1370}}
  >{\raggedright\arraybackslash}p{(\linewidth - 2\tabcolsep) * \real{0.8630}}@{}}
\toprule\noalign{}
\begin{minipage}[b]{\linewidth}\raggedright
Abbrev.
\end{minipage} & \begin{minipage}[b]{\linewidth}\raggedright
Expansion
\end{minipage} \\
\midrule\noalign{}
\endhead
\bottomrule\noalign{}
\endlastfoot
CI & Confidence interval (frequentist); credible interval (Bayesian) ---
context disambiguates. \\
RR & Relative risk (risk ratio). \\
OR & Odds ratio. \\
HR & Hazard ratio. \\
AUC & Area under the (ROC) curve. \\
LR+ / LR− & Positive / negative likelihood ratio. \\
PPV / NPV & Positive / negative predictive value. \\
MLE & Maximum likelihood estimator. \\
GLM & Generalised linear model. \\
GLMM & Generalised linear mixed model. \\
GEE & Generalised estimating equation. \\
RCT & Randomised controlled trial. \\
ITT & Intention to treat. \\
PP & Per protocol. \\
MI & Multiple imputation. \\
FDR & False discovery rate. \\
BH & Benjamini--Hochberg (FDR procedure). \\
TRIPOD & Transparent Reporting of a multivariable prediction model. \\
STROBE & Strengthening the Reporting of Observational Studies in
Epidemiology. \\
CONSORT & Consolidated Standards of Reporting Trials. \\
PRISMA & Preferred Reporting Items for Systematic Reviews and
Meta-Analyses. \\
ARRIVE & Animal Research: Reporting of In Vivo Experiments. \\
STARD & Standards for Reporting of Diagnostic Accuracy. \\
\end{longtable}
}

\subsection*{Code conventions}\label{code-conventions}
\addcontentsline{toc}{subsection}{Code conventions}

\begin{itemize}
\tightlist
\item
  The base R pipe \texttt{\textbar{}\textgreater{}} is used throughout.
  Files mixing \texttt{\%\textgreater{}\%} and
  \texttt{\textbar{}\textgreater{}} are flagged in
  \texttt{GENERATION\_LOG.md}.
\item
  Package functions are written qualified the first time they appear in
  a chapter (\texttt{dplyr::filter()}), then bare on subsequent uses.
\item
  Code blocks shown in the rendered book are display-only; the book is
  built with \texttt{eval\ =\ FALSE} globally and verified to knit
  cleanly without side effects. To run a chunk yourself, copy it into an
  R session with the libraries cited in that section attached.
\end{itemize}

\end{document}
